% LaTeX macros
\newcommand{\ghost}[1]{\raisebox{0pt}[0pt][0pt]{\makebox[0pt][l]{#1}}}
% math & logic
\newcommand{\N}{\mathbb{N}}
\newcommand{\Z}{\mathbb{Z}}
\newcommand{\R}{\mathbb{R}}
% number theory
%\newcommand{\divides}{\mid}
% why is this not a standard command but \gcd is?
%\newcommand{\lcm}{\mathbin{\mathrm{lcm}}}
% renew gcd to make it look like lcm
%\renewcommand{\gcd}{\mathbin{\mathrm{gcd}}}
% set theory
\newcommand{\set}[1]{\left\{#1\right\}}
\newcommand{\pset}[1]{2^{#1}}%
\newcommand{\sset}[1]{\left\{\,#1\,\right\}}
\newcommand{\guard}{\ \middle\vert\ }
\newcommand{\card}[1]{\left\vert#1\right\vert}
%\newcommand{\size}[1]{\left\Vert#1\right\Vert}
\newcommand{\tuple}[1]{\left(#1\right)}
%\newcommand{\tand}{\mathrel{\:\:\&\:\:}}
\newcommand{\tand}{\;\text{and}\;}
%\newcommand{\timplies}{\implies}
\newcommand{\timplies}{\quad\text{implies}\quad}
\newcommand{\define}[1]{\emph{#1}}
% use next command as define, but if you want something else (#2) to occur in the index/glossary
\newcommand{\defineas}[2]{\emph{#1}}
% thin spaces
\newcommand{\ts}{\mkern1mu}
\newcommand{\nts}{\mkern-1mu}
% ordinal numbers
\newcommand{\ordn}{\alpha}
\newcommand{\ordnalt}{\alpha'}
\newcommand{\ordsucc}{\mathord{\ordn+1}}
\newcommand{\ordlim}{\beta}
%% meta-level logic
% notation for meta-level quantification
%\newcommand{\qexists}[2]{(\exists#1)#2}
%\newcommand{\qforall}[2]{(\forall#1)#2}
%\newcommand{\onlyif}{\Longleftarrow}
% object-level macros
\newcommand{\eqdef}{\mathrel{:=}}
\newcommand{\iffdef}{\mathrel{\;:\!\!\iff\,}}
%% ebnf
% production rule/definition
\newcommand{\ebnfeq}{\mathrel{::=}}
% alternative
\newcommand{\ebnfalt}{\mathbin{\mid}}



% Define typographic struts, as suggested by Claudio Beccari
%   in an article in TeX and TUG News, Vol. 2, 1993.
\newcommand\Tstrut{\rule{0pt}{2.6ex}}       % = `top' strut
\newcommand\Bstrut{\rule[-1ex]{0pt}{0pt}}   % = `bottom' strut
\newcommand\Strut{\rule[-1ex]{0pt}{2.6ex}}  % = `both' strut

% Down right arrow
%\newcommand{\drarrow}{\rotatebox[origin=c]{180}{$\Lsh$}}







%%%%%%


%\newcommand{\citet}[1]{\citeauthor{#1}~[\citeyear{#1}]}

\newcommand{\explain}[1]{\ensuremath{#1}\dotfill\texttt{\detokenize{#1}}}

% explain macro and usage as table line
\newcommand{\taex}[2]{\ensuremath{#1}&\texttt{\detokenize{#1}}&#2\\}
\newcommand{\taexn}[2]{\textcolor{blue}{\ensuremath{#1}}&\textcolor{blue}{\texttt{\detokenize{#1}}}&\textcolor{blue}{#2}\\}
\newcommand{\taexp}[2]{\textcolor{purple}{\ensuremath{#1}}&\textcolor{purple}{\texttt{\detokenize{#1}}}&\textcolor{purple}{#2}\\}
\newcommand{\taexheading}[1]{\hline\multicolumn{3}{c}{\textbf{#1}}\\\hline}

\newlength{\narrowgatherskip}
\newenvironment{narrowgather}[1][2pt]{%
    \setlength{\narrowgatherskip}{#1}%
    \vskip\narrowgatherskip\mbox{}\hfill\(}{\)\hfill\mbox{}\vskip\narrowgatherskip\noindent}

